%========================================================================================
%  Seção 15.1 do D2L -- Word Embedding (word2vec)
%  https://d2l.ai/chapter_natural-language-processing-pretraining/word2vec.html
%========================================================================================

%------------------------------------------------
\section{Motivação: como representar palavras?}
%------------------------------------------------

\begin{frame}{Introdução: representação de palavras por vetores}
    \begin{center}
        {\large\color{dexpiBlue}\textbf{Se um computador só entende números como ele vai entender palavras?}}
    \end{center}

    \begin{columns}[c]
        \begin{column}{0.58\textwidth}
            \begin{itemize}
                \item O PLN resolve isso mapeando cada palavra para um \destaque{vetor de números reais} (\textbf{word embedding}).
                \item Palavras com \destaque{significados parecidos} ficam \destaque{próximas} matematicamente.
                \item Esses vetores são as \textit{features} de entrada que vão alimentar os nossos modelos.
            \end{itemize}
        \end{column}
        \begin{column}{0.38\textwidth}
            \centering
            \begin{tikzpicture}[font=\scriptsize, dot/.style={circle, fill=dexpiBlue, inner sep=1.3pt}]
                \draw[-{Stealth}, gray] (0,0) -- (3.6,0);
                \draw[-{Stealth}, gray] (0,0) -- (0,2.1);
                \draw[dexpiGreenDark, dashed, rounded corners] (0.3,1.05) rectangle (1.75,1.95);
                \draw[dexpiRedDark, dashed, rounded corners] (2.0,0.15) rectangle (3.45,1.0);
                \node[dot] at (0.6,1.6) {}; \node[right] at (0.6,1.6) {man};
                \node[dot] at (0.75,1.25) {}; \node[right] at (0.75,1.25) {son};
                \node[dot] at (2.3,0.75) {}; \node[right] at (2.3,0.75) {loves};
                \node[dot] at (2.45,0.4) {}; \node[right] at (2.45,0.4) {likes};
                \node[gray, anchor=north east] at (3.6,0) {$\mathbb{R}^d$};
            \end{tikzpicture}
        \end{column}
    \end{columns}

    \vspace{0.2cm}
    \centering
    \begin{tikzpicture}[
        box/.style={draw=dexpiBlue, fill=dexpiBlueLight, rounded corners=2pt, minimum height=0.8cm, align=center, font=\small},
        arr/.style={-{Stealth}, thick, dexpiBlue}]
        \node[box] (t) {Texto\\{\scriptsize ``the man loves his son''}};
        \node[box, right=0.35cm of t] (k) {Tokens\\{\scriptsize the | man | loves | \ldots}};
        \node[box, right=0.35cm of k] (i) {Índices\\{\scriptsize 0, 1, 2, 3, 4}};
        \node[box, right=0.35cm of i, fill=dexpiLightBlue, very thick] (v) {\textbf{Vetores}\\{\scriptsize $\xx_0, \xx_1, \ldots \in \mathbb{R}^d$}};
        \node[box, right=0.35cm of v] (m) {Modelo\\{\scriptsize RNN, Transformer\ldots}};
        \draw[arr] (t) -- (k);
        \draw[arr] (k) -- (i);
        \draw[arr] (i) -- (v);
        \draw[arr] (v) -- (m);
        \node[below=0.15cm of v, font=\scriptsize, text=dexpiRedDark] {\textbf{Tema de hoje}};
    \end{tikzpicture}
\end{frame}

\begin{frame}{Primeira tentativa: vetores \textit{one-hot}}
    \begin{columns}[c]
        \begin{column}{0.5\textwidth}
            \begin{itemize}
                \item Dicionário com $N = |\V|$ palavras, indexadas de $0$ a $N-1$.
                \item A palavra de índice $i$ vira um vetor de dimensão $N$ com \textbf{1} na posição $i$ e \textbf{0} no resto.
                \item É simples de construir e não exige treinamento.
            \end{itemize}
            \begin{block}{Exemplo ($N = 5$)}
                \small
                $\V = \{\text{the}, \text{man}, \text{loves}, \text{his}, \text{son}\}$
            \end{block}
        \end{column}
        \begin{column}{0.5\textwidth}
            \centering
            \begin{tikzpicture}[font=\small]
                \matrix (M) [matrix of math nodes, nodes={minimum width=0.6cm, minimum height=0.5cm, anchor=center},
                    left delimiter={[}, right delimiter={]}, row sep=1pt, column sep=1pt, ampersand replacement=\&] {
                    1 \& 0 \& 0 \& 0 \& 0 \\
                    0 \& 1 \& 0 \& 0 \& 0 \\
                    0 \& 0 \& 1 \& 0 \& 0 \\
                    0 \& 0 \& 0 \& 1 \& 0 \\
                    0 \& 0 \& 0 \& 0 \& 1 \\
                };
                \foreach \w [count=\r] in {the, man, loves, his, son} {
                    \node[anchor=east, text=dexpiBlue] at ($(M-\r-1.west)+(-0.35,0)$) {\textbf{\w}};
                }
                \foreach \c in {1,...,5} {
                    \fill[dexpiLightBlue, opacity=0.6] ($(M-\c-\c.north west)+(0.05,-0.02)$) rectangle ($(M-\c-\c.south east)+(-0.05,0.02)$);
                }
            \end{tikzpicture}
        \end{column}
    \end{columns}
\end{frame}

\begin{frame}{Por que \textit{one-hot} é uma má escolha?}
    \begin{itemize}
        \item Uma forma natural de medir semelhança entre vetores é a \destaque{similaridade do cosseno}:
    \end{itemize}
    \begin{equation*}
        \cos(\theta) = \frac{\xx^\top \yy}{\|\xx\|\,\|\yy\|} \in [-1, 1].
    \end{equation*}
    \begin{itemize}
        \item Para dois vetores one-hot de palavras \textbf{diferentes}: $\xx^\top\yy = 0 \Rightarrow \cos(\theta) = 0$.
        \item Ou seja: \textit{``man''} é tão parecido com \textit{``son''} quanto com \textit{``the''}. \textbf{Toda palavra é ortogonal a toda outra}.
    \end{itemize}

    \begin{columns}[t]
        \begin{column}{0.48\textwidth}
            \begin{alertblock}{Problemas}
                \small
                \begin{itemize}
                    \item Não codifica \textbf{semelhança} semântica.
                    \item Dimensão $N$ enorme ($10^5$--$10^6$) e vetores esparsos.
                \end{itemize}
            \end{alertblock}
        \end{column}
        \begin{column}{0.48\textwidth}
            \begin{exampleblock}{O que queremos}
                \small
                \begin{itemize}
                    \item Vetores \textbf{densos} de dimensão $d \ll N$ (ex.: $d = 100$--$300$).
                    \item Palavras parecidas $\Rightarrow$ vetores próximos.
                \end{itemize}
            \end{exampleblock}
        \end{column}
    \end{columns}
\end{frame}


\begin{frame}{Embedding: um dicionário de números}
    \begin{columns}[T]
        \begin{column}{0.5\textwidth}
            \small
            \begin{itemize}
                \item Pense num \destaque{dicionário}: você procura a palavra e recebe uma lista de números.
                \item Esse dicionário é uma tabela $\mathbf{W}$, com \textbf{uma linha por palavra}.
            \end{itemize}

            \begin{tikzpicture}
                \node[draw=dexpiBlue, fill=dexpiBlueLight, rounded corners=3pt, thick, inner sep=4pt, align=center] {
                    $\mathbf{e}_i^\top \mathbf{W} = \mathbf{W}_{i,:}$\\[1pt]
                    {\scriptsize multiplicar pelo one-hot $=$ \destaque{selecionar a linha da palavra}}};
            \end{tikzpicture}

            \vspace{0.3cm}
            \begin{tikzpicture}[baseline=-0.6ex]
                \foreach \a in {0,45,...,315} {
                    \fill[dexpiRedDark, rotate=\a] (-0.05,0.1) rectangle (0.05,0.18);
                }
                \fill[dexpiRedDark] (0,0) circle (0.13);
                \fill[white] (0,0) circle (0.05);
            \end{tikzpicture}
            {\footnotesize Os números são \textbf{\color{dexpiRedDark}APRENDIDOS}, não escolhidos por nós.}
        \end{column}
        \begin{column}{0.48\textwidth}
            \centering
            \begin{tikzpicture}
                \node[font=\small\bfseries, text=dexpiBlue] at (3.1,5.05) {O mapa do significado};
                \matrix (W) [matrix of nodes, ampersand replacement=\&, font=\scriptsize,
                    nodes={minimum height=0.34cm, anchor=center}, column sep=2pt,
                    column 1/.style={nodes={minimum width=1.0cm}},
                    column 2/.style={nodes={minimum width=0.75cm}},
                    column 3/.style={nodes={minimum width=0.75cm}}, row 2/.style={nodes={fill=dexpiGreenLight}}] at (3.1,3.95) {
                    rei \& 1.2 \& 1.5 \\
                    gato \& $-1.0$ \& 1.5 \\
                    correr \& $-0.6$ \& $-1.2$ \\
                };
                \draw[dexpiBlue] (W.north west) rectangle (W.south east);
                \draw[dexpiBlue] (W-1-1.north east) -- (W-3-1.south east);
                \node[font=\small, text=dexpiBlue, anchor=west] at ($(W.east)+(0.05,0)$) {$\mathbf{W}$};
                \begin{axis}[
                    at={(0,0)}, anchor=south west,
                    width=6.8cm, height=3.5cm, clip=false,
                    xmin=-2.4, xmax=2.6, ymin=-2.1, ymax=2.2,
                    axis lines=middle, ticks=none,
                    xlabel={\scriptsize dim 1}, ylabel={\scriptsize dim 2}]
                    \coordinate (gatoPt) at (axis cs:-1.0,1.5);
                    \addplot[only marks, mark=*, dexpiBlue] coordinates {(1.2,1.5) (1.6,1.1) (0.9,1.0)};
                    \node[font=\scriptsize, anchor=south] at (axis cs:1.2,1.5) {rei};
                    \node[font=\scriptsize, anchor=west] at (axis cs:1.6,1.1) {rainha};
                    \node[font=\scriptsize, anchor=east] at (axis cs:0.9,1.0) {príncipe};
                    \addplot[only marks, mark=square*, dexpiGreenDark] coordinates {(-1.4,1.2) (-1.0,1.5) (-1.6,0.8)};
                    \node[font=\scriptsize, anchor=south] at (axis cs:-1.0,1.5) {gato};
                    \node[font=\scriptsize, anchor=east] at (axis cs:-1.4,1.2) {cão};
                    \node[font=\scriptsize, anchor=north] at (axis cs:-1.6,0.8) {lobo};
                    \addplot[only marks, mark=triangle*, dexpiRed] coordinates {(-0.6,-1.2) (-1.1,-1.5) (-0.4,-1.6)};
                    \node[font=\scriptsize, anchor=south] at (axis cs:-0.6,-1.2) {correr};
                    \node[font=\scriptsize, anchor=east] at (axis cs:-1.1,-1.5) {andar};
                    \node[font=\scriptsize, anchor=west] at (axis cs:-0.4,-1.6) {pular};
                \end{axis}
                \draw[-{Stealth}, thick, dexpiGreenDark] (W-2-1.west) to[out=180, in=90] ($(gatoPt)+(0,0.4)$);
            \end{tikzpicture}\par
            {\scriptsize Cada linha de $\mathbf{W}$ é um ponto ($d=2$, ilustrativo).\\
            Significados parecidos $\Rightarrow$ \destaque{pontos próximos}.}
        \end{column}
    \end{columns}
\end{frame}

%------------------------------------------------
\section{word2vec: aprendizado auto-supervisionado}
%------------------------------------------------

\begin{frame}{word2vec \cite{mikolov2013efficient, mikolov2013distributed}}
    \begin{itemize}
        \item Proposto por Mikolov et al. (Google, 2013).
        \item Mapeia cada palavra para um vetor de \destaque{tamanho fixo} que expressa melhor \textbf{similaridade} e \textbf{analogia}.
        \item Dois modelos:
        \begin{itemize}
            \item \textbf{Skip-gram}: dada a palavra central, prevê as palavras do contexto.
            \item \textbf{CBOW} (\textit{Continuous Bag of Words}): dado o contexto, prevê a palavra central.
        \end{itemize}
    \end{itemize}

    \begin{block}{Por que \textit{auto-supervisionado}?}
        Não há rótulos humanos. Os ``rótulos'' (palavras vizinhas) saem do \textbf{próprio texto}. Qualquer corpus grande (Wikipédia, notícias, livros) vira dado de treino.
    \end{block}
\end{frame}

\begin{frame}{Janela de contexto: de onde vêm os dados de treino}
    \begin{itemize}
        \item Escolhemos uma \destaque{palavra central} $w_c$ e um \destaque{tamanho de janela} $m$.
        \item As palavras a até $m$ posições de distância formam o \destaque{contexto}.
    \end{itemize}
    \vspace{0.2cm}
    \centering
    \begin{tikzpicture}[
        w/.style={draw=dexpiBlue, rounded corners=2pt, minimum width=1.6cm, minimum height=0.8cm, font=\small},
        ctx/.style={w, fill=dexpiGreenLight, draw=dexpiGreenDark},
        cen/.style={w, fill=dexpiBlue, text=white, font=\small\bfseries}]
        \node[ctx] (a) at (0,0) {the};
        \node[ctx] (b) at (2,0) {man};
        \node[cen] (c) at (4,0) {loves};
        \node[ctx] (d) at (6,0) {his};
        \node[ctx] (e) at (8,0) {son};
        \draw[decorate, decoration={brace, amplitude=5pt, mirror}, thick, dexpiGreenDark]
            ($(a.south west)+(0,-0.1)$) -- ($(b.south east)+(0,-0.1)$) node[midway, below=6pt, font=\scriptsize] {$m=2$ à esquerda};
        \draw[decorate, decoration={brace, amplitude=5pt, mirror}, thick, dexpiGreenDark]
            ($(d.south west)+(0,-0.1)$) -- ($(e.south east)+(0,-0.1)$) node[midway, below=6pt, font=\scriptsize] {$m=2$ à direita};
        \node[above=0.1cm of c, font=\scriptsize, text=dexpiBlue] {centro $w^{(t)}$};
    \end{tikzpicture}

    \vspace{0.3cm}
    \begin{columns}[t]
        \begin{column}{0.45\textwidth}
            \small
            \textbf{Pares gerados (skip-gram):}\\
            (loves $\to$ the), (loves $\to$ man),\\
            (loves $\to$ his), (loves $\to$ son)
        \end{column}
        \begin{column}{0.45\textwidth}
            \small
            \textbf{Exemplo gerado (CBOW):}\\
            (\{the, man, his, son\} $\to$ loves)
        \end{column}
    \end{columns}
    \vspace{0.2cm}
    \raggedright
    \small Deslizando o centro por $t = 1, \ldots, T$, obtemos milhões de exemplos a partir de texto bruto.
\end{frame}

%------------------------------------------------
\section{O modelo Skip-gram}
%------------------------------------------------

\begin{frame}{diferenças entre Skip-gram e CBOW}
    \centering
    \begin{tikzpicture}[
        w/.style={draw=dexpiBlue, rounded corners=2pt, minimum width=1.4cm, minimum height=0.7cm, font=\small},
        ctx/.style={w, fill=dexpiGreenLight, draw=dexpiGreenDark},
        gap/.style={w, fill=white, dashed, very thick, font=\large\bfseries, text=dexpiRed},
        cen/.style={w, fill=dexpiBlue, text=white, font=\small\bfseries},
        lbl/.style={anchor=east, font=\small\bfseries, text=dexpiBlue, align=right}]
        % Skip-gram: centro -> vizinhos
        \node[lbl] at (-1.1,0) {Skip-gram\\[-2pt]{\scriptsize\mdseries adivinhar os vizinhos}};
        \node[gap] at (0,0) {?};
        \node[gap] at (1.7,0) {?};
        \node[cen] at (3.4,0) {loves};
        \node[gap] at (5.1,0) {?};
        \node[gap] at (6.8,0) {?};
        \node[font=\scriptsize, text=dexpiGreenDark, align=left, anchor=west] at (7.8,0) {resposta:\\the, man, his, son};
        % CBOW: contexto -> centro
        \node[lbl] at (-1.1,-1.2) {CBOW\\[-2pt]{\scriptsize\mdseries completar a lacuna}};
        \node[ctx] at (0,-1.2) {the};
        \node[ctx] at (1.7,-1.2) {man};
        \node[gap] at (3.4,-1.2) {?};
        \node[ctx] at (5.1,-1.2) {his};
        \node[ctx] at (6.8,-1.2) {son};
        \node[font=\scriptsize, text=dexpiGreenDark, align=left, anchor=west] at (7.8,-1.2) {resposta:\\loves};
    \end{tikzpicture}

    \vspace{0.2cm}
    \raggedright
    \small
    \begin{itemize}
        \item \textbf{Skip-gram}: ``se eu digo \textit{futebol}, o que aparece por perto?'' (\textit{bola, gol, time}\ldots). Cada vizinho vira uma \destaque{pergunta separada}: 4 vizinhos = 4 exemplos de treino.
        \item \textbf{CBOW}: o exercício de escola ``complete a frase''. Os vizinhos entram juntos, como um \destaque{saco de palavras} (\textit{bag of words}): a ordem não importa.
    \end{itemize}

    \begin{exampleblock}{Intuição (vale para os dois)}
        \small
        Para acertar, o modelo precisa aprender quais palavras combinam com quais contextos. Palavras com os \textbf{mesmos vizinhos} acabam com \textbf{vetores parecidos}.
    \end{exampleblock}
\end{frame}

\begin{frame}{Skip-gram passo a passo}
    \centering
    \begin{tikzpicture}[
        box/.style={draw=dexpiBlue, fill=dexpiBlueLight, rounded corners=2pt, align=center, font=\scriptsize, minimum height=0.8cm},
        num/.style={circle, fill=dexpiBlue, text=white, font=\scriptsize\bfseries, inner sep=1.5pt},
        arr/.style={-{Stealth}, thick, dexpiBlue}]
        \node[draw=dexpiBlue, fill=dexpiBlue, text=white, rounded corners=2pt, font=\scriptsize\bfseries, minimum width=1.2cm, minimum height=0.6cm] (c) at (0,0) {loves};
        \node[box, minimum width=1.6cm] (v) at (2.4,0) {vetor de\\centro $\vv$};
        \node[box, minimum width=2.4cm] (s) at (5.4,0) {compara com\\\textbf{cada palavra}\\do vocabulário\\(escores)};
        \node[box, fill=dexpiGreenLight, draw=dexpiGreenDark, minimum width=2.0cm] (p) at (8.6,0) {\textbf{softmax}\\probabilidades\\dos vizinhos};
        \node[box, fill=dexpiLightBlue, minimum width=2.6cm] (r) at (11.9,0) {confere com os\\vizinhos \textbf{reais}:\\the, man, his, son};
        \draw[arr] (c) -- (v);
        \draw[arr] (v) -- (s);
        \draw[arr] (s) -- (p);
        \draw[arr] (p) -- (r);
        \node[num] at ($(v.north)+(0,0.3)$) {1};
        \node[num] at ($(s.north)+(0,0.3)$) {2};
        \node[num] at ($(p.north)+(0,0.3)$) {3};
        \node[num] at ($(r.north)+(0,0.3)$) {4};
    \end{tikzpicture}

    \vspace{0.2cm}
    \raggedright
    \small
    \begin{enumerate}
        \item \textbf{Consulta}: pega o vetor de centro da palavra do meio.
        \item \textbf{Compara}: o quanto ele combina com cada candidata (produto interno).
        \item \textbf{Normaliza}: a softmax transforma os escores em probabilidades.
        \item \textbf{Confere}: quanto de probabilidade o modelo deu a cada vizinho verdadeiro?
    \end{enumerate}

    \begin{block}{Guarde para depois: o CBOW}
        \small
        Mesmo processo, mas \textbf{sem a média}: entra um vetor só e o palpite é conferido contra \textbf{vários} vizinhos.
    \end{block}
\end{frame}


\begin{frame}{A fórmula do skip-gram}
    \begin{equation*}
        P(\text{vizinho} \mid \text{centro}) =
        \frac{\exp\left({\color{dexpiBlue}\uu_o}^\top {\color{dexpiGreenDark}\vv_c}\right)}
             {{\color{dexpiRedDark}\sum_{i \in \V}} \exp\left(\uu_i^\top {\color{dexpiGreenDark}\vv_c}\right)}
    \end{equation*}

    \begin{columns}[t]
        \begin{column}{0.48\textwidth}
            \begin{block}{Lendo a fórmula}
                \small
                \begin{itemize}
                    \item {\color{dexpiGreenDark}$\vv_c$}: vetor de centro (passo 1).
                    \item {\color{dexpiBlue}$\uu_o$}: vetor da candidata a vizinha; o produto dá o escore (passo 2).
                    \item {\color{dexpiRedDark}$\sum$}: soma sobre todas as palavras = softmax (passo 3).
                \end{itemize}
            \end{block}
        \end{column}
        \begin{column}{0.48\textwidth}
            \begin{block}{Prévia do CBOW}
                \small
                O CBOW usará \textbf{esta mesma fórmula}, só trocando $\vv_c$ pela \textbf{média} dos vetores do contexto.
            \end{block}
        \end{column}
    \end{columns}

    \vspace{0.2cm}
    \small
    É uma classificação com $N$ classes: ``qual das $N$ palavras do vocabulário é a vizinha?''
\end{frame}

\begin{frame}{Como o modelo aprende: da probabilidade à perda}
    \small
    \begin{itemize}\setlength{\itemsep}{0pt}
        \item O modelo ajusta os vetores $\vv_i$ e $\uu_i$ de \textbf{todas} as palavras do vocabulário.
        \item A meta: dar \textbf{alta probabilidade} aos vizinhos que de fato aparecem no texto (\destaque{máxima verossimilhança}).
    \end{itemize}

    \vspace{0.1cm}
    {\centering
    \begin{tikzpicture}[
        passo/.style={draw=#1, fill=#1!10, line width=0.8pt, rounded corners=4pt,
                      text width=0.28\textwidth, minimum height=2.3cm, align=center, inner sep=3pt},
        seta/.style={-{Stealth[length=2.5mm]}, line width=1pt, dexpiBlueDark}]
        \node[passo=dexpiBlue] (p1) {%
            {\color{dexpiBlue}\large$\prod$\; \small\textbf{Passo 1: o produto}}\\[1pt]
            \footnotesize$\prod_{t}\prod_{j} P\big(w^{(t+j)} \mid w^{(t)}\big)$\\[1pt]
            \scriptsize Multiplica as probabilidades de \textbf{todos} os pares (centro $t$, vizinho $j$).\\[1pt]
            {\color{dexpiBlue}\scriptsize\textbf{maximizar}}};
        \node[passo=dexpiGreenDark, right=0.45cm of p1] (p2) {%
            {\color{dexpiGreenDark}\large$\log$\; \small\textbf{Passo 2: o log}}\\[1pt]
            \footnotesize$\sum_{t}\sum_{j} \log P\big(w^{(t+j)} \mid w^{(t)}\big)$\\[1pt]
            \scriptsize Transforma multiplicação em \textbf{soma} e evita números minúsculos.\\[1pt]
            {\color{dexpiGreenDark}\scriptsize\textbf{maximizar} (mesmo ótimo)}};
        \node[passo=dexpiRedDark, right=0.45cm of p2] (p3) {%
            {\color{dexpiRedDark}\large$-$\; \small\textbf{Passo 3: o negativo}}\\[1pt]
            \footnotesize$\mathcal{L} = -\sum_{t}\sum_{j} \log P\big(w^{(t+j)} \mid w^{(t)}\big)$\\[1pt]
            \scriptsize ``Maximizar o acerto'' vira ``\textbf{minimizar o erro}'' (= entropia cruzada).\\[1pt]
            {\color{dexpiRedDark}\scriptsize\textbf{minimizar}}};
        \draw[seta] (p1) -- (p2);
        \draw[seta] (p2) -- (p3);
    \end{tikzpicture}\par}

    \vspace{0.1cm}
    \begin{exampleblock}{O que é $P\big(w^{(t+j)} \mid w^{(t)}\big)$?}
        \footnotesize
        É a \textbf{similaridade} entre os vetores ($\uu_o^\top \vv_c$) passada pelo \destaque{softmax}: a fórmula do slide anterior.
    \end{exampleblock}
\end{frame}

\begin{frame}{Otimizando com SGD: um par de palavras por vez}
    \small
    A perda mede o erro; falta saber \destaque{para que lado} empurrar cada vetor (o \textbf{gradiente}). Com bilhões de palavras, cada passo usa um \textbf{trecho curto}: \destaque{SGD}.

    \vspace{0.15cm}
    {\centering
    \begin{tikzpicture}[
        carta/.style={draw=#1, fill=#1!8, line width=0.8pt, rounded corners=4pt,
                      text width=0.27\textwidth, minimum height=1.55cm, align=center, inner sep=3pt},
        seta/.style={-{Stealth[length=2.5mm]}, line width=1pt, dexpiBlueDark}]
        \node[carta=dexpiBlue] (c1) {%
            {\color{dexpiBlue}\large\textbf{1}}\;\; \small\textbf{um trecho}\\[1pt]
            {\footnotesize cada passo olha}\\[0pt]
            {\footnotesize um \textbf{lote} pequeno}\\[1pt]
            {\color{dexpiBlue}\scriptsize estocástico $=$ sorteado}};
        \node[carta=dexpiGreenDark, right=0.45cm of c1] (c2) {%
            {\color{dexpiGreenDark}\large\textbf{2}}\;\; \small\textbf{um par}\\[1pt]
            {\footnotesize bastam \textbf{dois} vetores:}\\[0pt]
            {\footnotesize centro e vizinho}\\[1pt]
            {\color{dexpiGreenDark}\scriptsize $\log P(w_o \mid w_c)$}};
        \node[carta=dexpiRedDark, right=0.45cm of c2] (c3) {%
            {\color{dexpiRedDark}\large\textbf{3}}\;\; \small\textbf{um passo}\\[1pt]
            {\footnotesize anda no sentido}\\[0pt]
            {\footnotesize que \textbf{diminui} a perda}\\[1pt]
            {\color{dexpiRedDark}\scriptsize $\vv \leftarrow \vv - \eta\,\nabla$}};
        \draw[seta] (c1) -- (c2);
        \draw[seta] (c2) -- (c3);
    \end{tikzpicture}\par}

    \vspace{0.2cm}
    {\centering
    $\log P(w_o \mid w_c) =
     \underbrace{{\color{dexpiBlue}\uu_o^\top \vv_c}}_{\text{nota do vizinho real}}
     \;-\;
     \underbrace{{\color{dexpiRedDark}\log\Big(\sum_{i \in \V} \exp(\uu_i^\top \vv_c)\Big)}}_{\text{soma sobre o vocabulário}}$
    \par}

    \vspace{0.15cm}
    \begin{exampleblock}{Uma metade premia, a outra pune}
        \footnotesize
        A {\color{dexpiBlue}nota do vizinho real} quer subir; a {\color{dexpiRedDark}soma} pune qualquer nota alta, inclusive a das palavras erradas.
    \end{exampleblock}
\end{frame}

\begin{frame}{O gradiente: observado menos esperado}
    \footnotesize
    Derivar o termo do slide anterior usa só duas regras que já conhecemos:
    \quad
    $\dfrac{\partial}{\partial \vv_c}\big(\uu_o^\top \vv_c\big) = \uu_o$
    \quad e \quad
    $\dfrac{\partial}{\partial \vv_c}\log\sum_{i \in \V}\exp(\uu_i^\top \vv_c) = \sum_{j \in \V} P(w_j \mid w_c)\,\uu_j$.

    \vspace{0.15cm}
    {\centering
    $\nabla_{\vv_c}\log P(w_o \mid w_c) =
     \underbrace{{\color{dexpiBlue}\uu_o}}_{\text{observado}}
     \;-\;
     \underbrace{{\color{dexpiRedDark}\sum_{j \in \V} P(w_j \mid w_c)\,\uu_j}}_{\text{esperado pelo modelo}}$
    \par}

    \vspace{0.1cm}
    \begin{columns}[T]
        \begin{column}{0.53\textwidth}
            \begin{block}{Lendo o resultado}
                \footnotesize
                \begin{itemize}\setlength{\itemsep}{1pt}
                    \item Se o modelo já acerta, esperado $\approx$ observado e $\nabla \approx 0$: não há o que corrigir.
                    \item A soma percorre as $N$ palavras do vocabulário, a cada par: é o custo que a Seção 15.2 resolve.
                \end{itemize}
            \end{block}
        \end{column}
        \begin{column}{0.43\textwidth}
            \centering
            \begin{tikzpicture}[scale=1.75, font=\scriptsize]
                \draw[gray!55, -{Stealth}] (-1.2,0) -- (1.3,0);
                \draw[gray!55, -{Stealth}] (0,-0.45) -- (0,1.2);
                \node[circle, fill=gray!55, inner sep=1.6pt] (man) at (0,1) {};
                \node[left=3pt of man] {man};
                \node[circle, fill=gray!55, inner sep=1.6pt] (the) at (-1,0) {};
                \node[above=2pt of the] {the};
                \node[circle, fill=dexpiBlue, inner sep=1.7pt] (son) at (1,0) {};
                \node[below right=1pt and 1pt of son, text=dexpiBlue] {son};
                \node[circle, draw=dexpiBlue, fill=white, inner sep=1.4pt] (esp) at (0.56,0.26) {};
                \node[above left=2pt and 0pt of esp, text=dexpiBlue] {esperado};
                \draw[-{Stealth}, very thick, dexpiBlue] (esp) -- (son);
                \node[above right=0pt and 3pt of son, text=dexpiBlue] {$\nabla$};
            \end{tikzpicture}\par
            \vspace{0.05cm}
            {\scriptsize A seta vai do que o modelo \textbf{esperava} até o vizinho que \textbf{apareceu}.}
        \end{column}
    \end{columns}
\end{frame}

%------------------------------------------------
\section{O modelo CBOW}
%------------------------------------------------

\begin{frame}{CBOW passo a passo}
    \centering
    \begin{tikzpicture}[
        ctx/.style={draw=dexpiGreenDark, fill=dexpiGreenLight, rounded corners=2pt, minimum width=1.2cm, minimum height=0.5cm, font=\scriptsize},
        box/.style={draw=dexpiBlue, fill=dexpiBlueLight, rounded corners=2pt, align=center, font=\scriptsize, minimum height=0.8cm},
        num/.style={circle, fill=dexpiBlue, text=white, font=\scriptsize\bfseries, inner sep=1.5pt},
        arr/.style={-{Stealth}, thick, dexpiBlue}]
        % passo 1
        \foreach \t [count=\i from 0] in {the, man, his, son} {
            \node[ctx] (c\i) at (0, 0.9-\i*0.6) {\t};
            \node[box, minimum height=0.5cm, right=0.4cm of c\i] (v\i) {vetor};
            \draw[arr] (c\i) -- (v\i);
        }
        % passo 2
        \node[box, fill=dexpiLightBlue, minimum width=2.0cm] (m) at (5.2,0) {\textbf{média}\\``resumo do\\contexto''};
        \foreach \i in {0,...,3} { \draw[arr] (v\i.east) -- (m.west); }
        % passo 3
        \node[box, minimum width=2.4cm] (s) at (8.6,0) {compara com\\\textbf{cada palavra}\\do vocabulário\\(escores)};
        \draw[arr] (m) -- (s);
        % passo 4
        \node[box, fill=dexpiGreenLight, draw=dexpiGreenDark, minimum width=2.2cm] (p) at (12.0,0) {\textbf{softmax}\\probabilidades\\$\Rightarrow$ ``loves''};
        \draw[arr] (s) -- (p);
        % números
        \node[num] at (1.4,1.45) {1};
        \node[num] at ($(m.north)+(0,0.3)$) {2};
        \node[num] at ($(s.north)+(0,0.3)$) {3};
        \node[num] at ($(p.north)+(0,0.3)$) {4};
    \end{tikzpicture}

    \vspace{0.2cm}
    \raggedright
    \small
    \begin{enumerate}
        \item \textbf{Consulta}: pega o vetor de cada palavra do contexto na tabela de embeddings.
        \item \textbf{Resume}: tira a média, um único vetor que ``representa o contexto''.
        \item \textbf{Compara}: o quanto o resumo combina com cada candidata (produto interno).
        \item \textbf{Normaliza}: a softmax transforma os escores em probabilidades.
    \end{enumerate}
\end{frame}


\begin{frame}{A fórmula do CBOW}
    \begin{equation*}
        P(\text{centro} \mid \text{contexto}) =
        \frac{\exp\left({\color{dexpiBlue}\uu_c}^\top {\color{dexpiGreenDark}\bar{\vv}_o}\right)}
             {{\color{dexpiRedDark}\sum_{i \in \V}} \exp\left(\uu_i^\top {\color{dexpiGreenDark}\bar{\vv}_o}\right)}
        \qquad \text{com} \quad {\color{dexpiGreenDark}\bar{\vv}_o} = \frac{\vv_{o_1} + \cdots + \vv_{o_{2m}}}{2m}
    \end{equation*}

    \begin{columns}[t]
        \begin{column}{0.48\textwidth}
            \begin{block}{Lendo a fórmula}
                \small
                \begin{itemize}
                    \item {\color{dexpiGreenDark}$\bar{\vv}_o$}: o resumo do contexto (passo 2).
                    \item {\color{dexpiBlue}$\uu_c$}: o vetor da candidata; o produto dá o escore (passo 3).
                    \item {\color{dexpiRedDark}$\sum$}: soma sobre todas as palavras = softmax (passo 4).
                \end{itemize}
            \end{block}
        \end{column}
        \begin{column}{0.48\textwidth}
            \begin{block}{Relação com o skip-gram}
                \small
                \textbf{É a mesma fórmula!} Só troca o vetor do centro pela \textbf{média} dos vetores do contexto.
            \end{block}
        \end{column}
    \end{columns}

    \vspace{0.1cm}
    \begin{exampleblock}{Cuidado com a notação do D2L}
        \footnotesize
        No CBOW o livro inverte as letras: $\vv$ são os vetores das palavras de \textbf{contexto} (a entrada) e $\uu_c$ é o vetor do \textbf{centro} previsto. Regra para não se perder: $\vv$ é sempre o vetor de quem \textbf{entra} no modelo.
    \end{exampleblock}
\end{frame}

\begin{frame}{Treinando o CBOW: a mesma perda, um termo por posição}
    \small
    \begin{itemize}\setlength{\itemsep}{1pt}
        \item A mesma receita de antes: \destaque{máxima verossimilhança}, $\log$ e sinal de menos.
        \item O que muda é \textbf{quantos termos} cada posição do texto gera.
    \end{itemize}

    \vspace{0.1cm}
    \begin{columns}[c]
        \begin{column}{0.48\textwidth}
            \begin{block}{Skip-gram: $2m$ termos}
                \footnotesize
                $\mathcal{L} = -\sum_{t}\sum_{j \neq 0} \log P(w^{(t+j)}\mid w^{(t)})$
                \\[2pt]
                Um termo para \textbf{cada vizinho} da janela ($-m \le j \le m$, $j \neq 0$): com $m = 2$, são \destaque{4} termos.
            \end{block}
        \end{column}
        \begin{column}{0.48\textwidth}
            \begin{block}{CBOW: \textbf{1} termo}
                \footnotesize
                $\mathcal{L} = -\sum_{t} \log P(w^{(t)}\mid C^{(t)})$
                \\[2pt]
                O contexto $C^{(t)}$ entra todo \textbf{junto}, como um saco de palavras: com $m = 2$, é \destaque{1} termo.
            \end{block}
        \end{column}
    \end{columns}

    \vspace{0.15cm}
    \begin{exampleblock}{O que se ganha e o que se perde}
        \footnotesize
        \textbf{Ganha}: menos termos por posição $\Rightarrow$ treino mais rápido. \quad
        \textbf{Perde}: a média ignora a ordem das palavras --- é o preço do ``saco de palavras''.
    \end{exampleblock}
\end{frame}

\begin{frame}{O gradiente do CBOW: o mesmo, dividido por $2m$}
    \small
    \begin{itemize}\setlength{\itemsep}{1pt}
        \item O log do CBOW é o do skip-gram trocando $\vv_c$ por $\bar{\vv}_o$; então a derivada também é a mesma. A novidade é derivar a \textbf{média} (aqui $C$ é o contexto da posição).
    \end{itemize}
    \vspace{0.05cm}
    {\centering
    $\frac{\partial \log P(w_c \mid C)}{\partial \vv_{o_i}} =
     \frac{1}{2m}\Big(
     \underbrace{{\color{dexpiBlue}\uu_c}}_{\text{observado}} -
     \underbrace{{\color{dexpiRedDark}\sum_{j \in \V} P(w_j \mid C)\,\uu_j}}_{\text{esperado pelo modelo}}\Big)$
    \par}

    \vspace{0.2cm}
    {\centering
    \begin{tikzpicture}[
        caixa/.style={draw=#1, fill=#1!8, rounded corners=3pt, line width=0.8pt,
                      text width=0.26\textwidth, minimum height=1.45cm, align=center, inner sep=3pt},
        seta/.style={-{Stealth[length=2.5mm]}, line width=1pt, dexpiBlueDark}]
        \node[caixa=dexpiBlue] (b1) {%
            \footnotesize$\bar{\vv}_o = \frac{1}{2m}\sum_{i=1}^{2m}\vv_{o_i}$\\[2pt]
            {\scriptsize resumo do contexto}};
        \node[caixa=dexpiGreenDark, right=0.45cm of b1] (b2) {%
            \footnotesize$\frac{\partial \bar{\vv}_o}{\partial \vv_{o_i}} = \frac{1}{2m}$\\[2pt]
            {\scriptsize cada palavra pesa $1/2m$}};
        \node[caixa=dexpiRedDark, right=0.45cm of b2] (b3) {%
            \footnotesize o ajuste chega \textbf{repartido}\\[2pt]
            {\scriptsize entre as $2m$ palavras do contexto}};
        \draw[seta] (b1) -- (b2);
        \draw[seta] (b2) -- (b3);
    \end{tikzpicture}\par}

    \vspace{0.15cm}
    \begin{exampleblock}{O que isso muda na prática}
        \footnotesize
        Todas as palavras do contexto são puxadas \textbf{na mesma direção}, com um passo menor. É por isso que o CBOW dilui palavras raras: elas aparecem poucas vezes e sempre dentro de uma média.
    \end{exampleblock}
\end{frame}

%------------------------------------------------
\section{Implementação e resumo}
%------------------------------------------------

\begin{frame}[fragile]{Implementação mínima em PyTorch (softmax completa)}
\begin{columns}[T]
    \begin{column}{0.62\textwidth}
\begin{lstlisting}[basicstyle=\ttfamily\fontsize{7}{8}\selectfont,aboveskip=0pt,belowskip=0pt]
class SkipGram(nn.Module):
    def __init__(self, vocab, dim):
        super().__init__()
        self.V = nn.Embedding(vocab, dim)  # centro
        self.U = nn.Embedding(vocab, dim)  # contexto

    def forward(self, center):
        v = self.V(center)          # (B, d)
        return v @ self.U.weight.T  # (B, N)

model = SkipGram(10_000, 100)
opt = torch.optim.SGD(model.parameters(), 0.01)

logits = model(center)
loss = F.cross_entropy(logits, context)
loss.backward(); opt.step()
emb = model.V.weight            # vetor final
\end{lstlisting}
    \end{column}
    \begin{column}{0.34\textwidth}
        \footnotesize
        \textbf{Onde está cada passo}\\[2pt]
        \begin{enumerate}\setlength{\itemsep}{2pt}
            \item \textbf{Consulta}: \texttt{self.V(center)}.
            \item \textbf{Compara}: \texttt{v @ U.weight.T} mede a afinidade com as $N$ palavras.
            \item \textbf{Normaliza}: a \texttt{cross\_entropy} já aplica a softmax e o $-\log$.
            \item \textbf{Confere}: o alvo \texttt{context} é o vizinho real.
        \end{enumerate}
        \vspace{0.1cm}
        {\scriptsize\color{dexpiRedDark} \textbf{O gargalo:} $N = 10^4$ já são $10^4$ contas por par. É o que a próxima seção resolve.}
    \end{column}
\end{columns}
\end{frame}
